% ===== PRÉAMBULE CANONIQUE — à coller VERBATIM en tête de CHAQUE .tex =====
\documentclass[11pt,a4paper]{article}
\usepackage[utf8]{inputenc}\usepackage[T1]{fontenc}\usepackage{lmodern}
\usepackage[french]{babel}
\usepackage{amsmath,amssymb,mathtools}\usepackage{icomma}
\usepackage[version=4]{mhchem}          % \ce{...} : équations & formules
\usepackage{siunitx}\sisetup{locale=FR} % \qty{}{}, \si{}, \ang{}
\usepackage{chemfig}                    % \chemfig{...} : structures développées
\usepackage{xcolor}\usepackage{colortbl}
\usepackage{tikz}\usepackage{pgfplots}\pgfplotsset{compat=1.18}
\usepackage[most]{tcolorbox}\usepackage{enumitem}\usepackage{array,booktabs}
\usepackage[a4paper,margin=2.2cm]{geometry}
\usetikzlibrary{arrows.meta,calc,angles,quotes,positioning,patterns,backgrounds,shapes.geometric}
\definecolor{primary}{RGB}{222,49,99}\definecolor{primarydark}{RGB}{150,22,55}
\definecolor{defcol}{RGB}{0,114,178}\definecolor{methcol}{RGB}{52,92,148}
\definecolor{excol}{RGB}{0,135,90}\definecolor{attncol}{RGB}{200,40,55}
\tcbset{boxbase/.style={enhanced,breakable,boxrule=0.9pt,arc=2.5pt,
  left=3mm,right=3mm,top=2.3mm,bottom=2.3mm,fonttitle=\bfseries,coltitle=white}}
% --- boîtes TUTORIEL (noms EXACTS, ne pas renommer) ---
\newtcolorbox{cle}[1][]{boxbase,colback=defcol!6,colframe=defcol,
  title={\textsc{À retenir}\ifx\relax#1\relax\relax\else\textmd{ : \itshape #1}\fi}}
\newtcolorbox{methode}[1][]{boxbase,colback=methcol!7,colframe=methcol,sharp corners,
  title={\textsc{Méthode}\ifx\relax#1\relax\relax\else\textmd{ --- \itshape #1}\fi}}
\newtcolorbox{exemplebox}[1][]{boxbase,colback=excol!5,colframe=excol,
  title={\textsc{Exemple}\ifx\relax#1\relax\relax\else\textmd{ : \itshape #1}\fi}}
\newtcolorbox{piege}[1][]{boxbase,colback=attncol!6,colframe=attncol,
  title={\textsc{Piège}\ifx\relax#1\relax\relax\else\textmd{ : \itshape #1}\fi}}
\newtcolorbox{remarque}[1][]{boxbase,colback=black!4,colframe=black!45,
  title={\textsc{Remarque}\ifx\relax#1\relax\relax\else\textmd{ : \itshape #1}\fi}}
% --- boîtes EXERCICE / CORRIGÉ (noms EXACTS, auto-comptées) ---
\newtcolorbox[auto counter]{exo}[1][]{boxbase,colback=primary!4,colframe=primary,
  title={\textsc{Exercice}~\thetcbcounter\ifx\relax#1\relax\relax\else\textmd{ --- \itshape #1}\fi}}
\newtcolorbox[auto counter]{corrige}[1][]{boxbase,colback=excol!4,colframe=excol,
  title={\textsc{Corrigé}~\thetcbcounter\ifx\relax#1\relax\relax\else\textmd{ --- \itshape #1}\fi}}
\newcommand{\R}{\mathbb{R}}\newcommand{\N}{\mathbb{N}}\newcommand{\Z}{\mathbb{Z}}\newcommand{\ds}{\displaystyle}
% (un \usepackage{fancyhdr}+en-tête est possible ; il est ignoré côté web)
% =========================================================================
\begin{document}

\begin{corrige}[nommer un oxyde métallique à valence fixe]
Métaux à valence unique : pas de chiffre romain.
\begin{enumerate}[label=\alph*)]
  \item \ce{CaO} : oxyde de calcium.
  \item \ce{K2O} : oxyde de potassium.
  \item \ce{ZnO} : oxyde de zinc.
  \item \ce{Al2O3} : oxyde d'aluminium.
\end{enumerate}
\end{corrige}

\begin{corrige}[écrire un oxyde métallique à valence fixe]
On croise la valence du métal avec $2$ (celle de \ce{O^2-}), puis on simplifie.
\begin{enumerate}[label=\alph*)]
  \item oxyde de sodium : \ce{Na+} (valence 1) $\Rightarrow$ \ce{Na2O}.
  \item oxyde de magnésium : \ce{Mg^2+} (valence 2) $\Rightarrow$ \ce{Mg2O2}\,$\to$\,\ce{MgO}.
  \item oxyde d'argent : \ce{Ag+} (valence 1) $\Rightarrow$ \ce{Ag2O}.
\end{enumerate}
\end{corrige}

\begin{corrige}[nommer un oxyde non métallique \ce{XO_n}]
Un seul atome de $X$ : le préfixe est le nombre d'oxygènes.
\begin{enumerate}[label=\alph*)]
  \item \ce{CO} : monoxyde de carbone.
  \item \ce{CO2} : dioxyde de carbone.
  \item \ce{SO3} : trioxyde de soufre.
  \item \ce{NO2} : dioxyde d'azote.
\end{enumerate}
\end{corrige}

\begin{corrige}[écrire un oxyde non métallique à partir du préfixe]
\begin{enumerate}[label=\alph*)]
  \item monoxyde d'azote : \ce{NO}.
  \item dioxyde d'iode : \ce{IO2}.
  \item trioxyde de brome : \ce{BrO3}.
\end{enumerate}
\end{corrige}

\begin{corrige}[choisir la bonne règle : métal ou non-métal ?]
\begin{enumerate}[label=\alph*)]
  \item \ce{MgO} : oxyde \emph{métallique} (magnésium) $\Rightarrow$ oxyde de magnésium.
  \item \ce{SO2} : oxyde \emph{non métallique} (soufre) $\Rightarrow$ dioxyde de soufre.
  \item \ce{Na2O} : oxyde \emph{métallique} (sodium) $\Rightarrow$ oxyde de sodium.
\end{enumerate}
\end{corrige}

\end{document}
